\newpage
\section{Preliminaries}
\label{sec:prelims}

\subsection{Matrix notation}

Given matrices
$\bA, \bB \in \R^{n\times n}$, we will use:
\begin{itemize}
    \item $\bA \in \R^{n\times n}_\sym$ to
    specify that $\bA$ is symmetric.
    \item $\bA \in O(n) \subseteq \R^{n \times n}$ to specify that $\bA$ is orthogonal.
    \item $\bA[i,j]$ to denote its $(i,j)$-th
    entry for $i,j\in [n] \defeq \{1,\ldots,n\}$.
    \item $\|\bA\| \defeq \max_{\|\bm x\|_2=1} \|\bA \bm x\|_2$ to denote its spectral or operator
    norm.
    \item $\|\bA\|^2_\frob \defeq \sum_{i,j=1}^n \bA[i,j]^2$ to denote its
    Frobenius norm.
    \item $\Tr(\bm A) \defeq \sum_{i=1}^n \bm A[i,i]$ to denote its trace.
    \item $\lambda_1(\bA)\ge \ldots \ge \lambda_n(\bA)$ to
    denote its eigenvalues when $\bm A$ is symmetric.
    \item $\bA \odot \bB$ to denote the entrywise or Hadamard product with entries $(\bm A[i,j]\bm B[i,j])_{i,j\in [n]}$.
\end{itemize}

\begin{definition}[Puncturing]\label{def:puncturing}
    Let $\bH\in \R_\sym^{n\times n}$ and $\bm \Pi \defeq \bm I - \frac 1n \bm 1 \bm 1^\top$ be
    the projection orthogonal to the all-ones
    direction. The {puncturing} of $\bH$ is
    the matrix $\bA = \bm \Pi \bmH \bm \Pi$.
\end{definition}

\begin{definition}[\goe]
\label{def:goe}
    The (normalized) {Gaussian Orthogonal Ensemble} \goe
    is the distribution of random matrices $\bA\in\R_\sym^{n\times n}$
    with $\bA[i,j]=\bA[j,i]\sim \calN(0,1/n)$ independently for all $1\le i<j\le n$, and
    $\bA[i, i]\sim \calN(0,2/n)$ independently
    for all $i\in [n]$.
\end{definition}

\begin{definition}[Hadamard matrices]
    \label{def:hadamard}
    When $n$ is a power of $2$, the (normalized) Walsh--Hadamard matrix $\bH_{\textnormal{had}}^{(n)}\in\R_\sym^{n\times n}$ is defined 
    recursively by
    \[
        \bH_{\textnormal{had}}^{(1)}=\begin{bmatrix}1\end{bmatrix}\,,\qquad 
        \bH_{\textnormal{had}}^{(2n)} \defeq \frac{1}{\sqrt{2}}
        \begin{bmatrix}
        \bH_{\textnormal{had}}^{(n)} & \bH_{\textnormal{had}}^{(n)}\\
        \bH_{\textnormal{had}}^{(n)} & -\bH_{\textnormal{had}}^{(n)}
        \end{bmatrix}.
    \]
\end{definition}
\noindent
$\bH_{\textnormal{had}}^{(n)}$ is a 
symmetric orthogonal matrix with entries in $\pm 1/\sqrt n$.


\begin{definition}[DST and DCT matrices]
    \label{def:dst}
    The discrete sine transform matrices
    $\bH_{\sin}^{(n)}\in\R^{n\times n}_\sym$ are
    \[
        \bH_{\sin}^{(n)}[i,j] \defeq \sqrt{\frac 2 {n+1}} \sin\left(\frac {\pi ij}{n+1}\right)\quad \forall i,j\in [n]\,.
    \]
    The discrete cosine transform matrices
    $\bH_{\cos}^{(n)}\in\R^{n\times n}_\sym$ are 
    \[
        \bH_{\cos}^{(n)}[i,j]\defeq \sqrt{\frac 2n}\cos\left(\frac{\pi(i-\tfrac12)(j-\tfrac12)}{n}\right)\quad \forall i,j\in [n]\,.
    \]
\end{definition}
\noindent
$\bH_{\cos}^{(n)}$ and $\bH_{\sin}^{(n)}$ are
symmetric orthogonal matrices with entries
at most $O(1/\sqrt n)$ in magnitude.

\begin{definition}[\rom and \rrom]
\label{def:rom}
    The {Random Orthogonal Model} \rom is the
    distribution of random matrices $\bH = \bm Q \bm D \bm Q^\top$, where $\bm Q\in O(n)$ is
    Haar-distributed, and $\bm D$ is a diagonal
    matrix with i.i.d. $\textnormal{Unif}(\{-1,1\})$ entries, independent from $\bm Q$.
    The {Regular Random Orthogonal Model} \rrom is the distribution of the puncturing of
    $\bH$, when $\bH$ is sampled from the \rom.
\end{definition}
\noindent
Random matrices from the \rom are symmetric orthogonal matrices, satisfying $\bH^2 = \bm I$.
They are a special case of the
orthogonally invariant models we discuss in~\Cref{sec:rot-inv}.

\subsection{Modes of convergence}
\label{sec:prelim-modes}

We will use a few standard modes of convergence from scalar-valued probability theory.
\begin{definition}[Modes of convergence: scalars]
    For a sequence of random variables $x^{(n)}\in \R$, we say that:
    \begin{itemize}
        \item $x^{(n)}$ \emph{converge in expectation} if, for some $c \in \R$, $\lim_{n \to \infty} \EE x^{(n)} = c$.
        \item $x^{(n)}$ \emph{converge in probability} if, for some $c \in \R$, for all $\varepsilon > 0$, $\lim_{n \to \infty} \PP[|x^{(n)} - c| > \varepsilon] = 0$.
        \item $x^{(n)}$ \emph{converge in $L^2$} if they converge in expectation and $\lim_{n \to \infty} \EE(x^{(n)} - c)^2 = 0$, or equivalently if they converge in expectation and $\lim_{n \to \infty} \Var x^{(n)} = 0$.
    \end{itemize}
    We write a symbol $\cM \in \{\EE, \PP, L^2\}$ to indicate these modes of convergence, and in this notation say that the $x^{(n)}$ \emph{converge in $\cM$}.
\end{definition}

Moreover, we say a sequence of random vectors
$\bx^{(n)}\in\R^d$ in fixed dimension
$d\ge 1$ {\em converges in distribution} to a random vector
$\bx\in\R^d$ if for every bounded continuous
function $\varphi \colon \R^d\to \R$,
\[
    \E \varphi(\bx^{(n)})\underset{n\to\infty}\longrightarrow \E \varphi(\bx)\,,
\]
in which case we write $\bx^{(n)}\tod \bx$.
See~\Cref{sec:convergence-process} for a
generalization to random variables
indexed by a countably infinite index
set.

\begin{definition}[Modes of convergence: tracial moments]
    For a mode of convergence $\cM$, we say that a sequence of random matrices $\bA \in \R^{n \times n}$ \emph{converges in tracial moments in $\cM$} if, for every $k \geq 1$, $\frac{1}{n}\Tr \bA^k$ converges in $\cM$.
    We say that it converges in tracial moments in $\cM$ to a probability measure $\mu$ over $\R$ if \[\frac{1}{n}\Tr \bA^k \to \int x^k\, \d\mu(x)\]in the mode of convergence $\cM$.
\end{definition}

\subsection{Matchings and Wick calculus}

Given a set $S$, let $\calM(S)$ denote the set of matchings on $S$.
Let $\PM(S)$ denote the subset of perfect matchings.
The elements of $M \in \calM(S)$ are written as pairs $\{i,j\} \subseteq S$.
For several sets $S_1, \dots, S_k$, denote by $\calM(S_1, \dots, S_k)$ the set of matchings on the disjoint union $S_1 \sqcup \cdots \sqcup S_k$ that do not match any two elements of the same $S_i$.
For two sets $S_1, S_2$ of the same size, denote by $\PM(S_1, S_2)$ the bipartite perfect matchings of $S_1 \sqcup S_2$ that only match elements of $S_1$ to ones of $S_2$.
We will abbreviate $\calM(\{1,2,\dots, q\})$ as $\calM(q)$.

\begin{lemma}[Wick lemma]\label{lem:wick}
    Let $X_1, \dots, X_q$ be jointly Gaussian random variables with mean zero. Then:
    \begin{equation*}
        \E[X_1\cdots X_{q}] = \sum_{M \in \PM(q)} \prod_{ij \in M} \E[X_iX_j]\,.
    \end{equation*}
\end{lemma}

The Wick products are the multivariate generalization of the Hermite polynomials to correlated Gaussians~\cite[Chapter 3]{Janson:GaussianHilbertSpaces}.

\begin{definition}[Wick product]
\label{def:wick}
    Let $I$ be an index set, $\bX = (X_i)_{i \in I}$ be formal variables, and $\matSig \in \R_{\sym}^{I \times I}$. The {\em Wick products} are defined by, for each finitely supported $\alpha \in \N^I$,
    \[
       \He_{\al}(\bX\,;\, \matSig) \defeq \sum_{M \in \calM(\al)} (-1)^{|M|} \prod_{uv \in M}\matSig[u,v] \prod_{u\notin M} X_u\,,
    \]
    where $\calM(\al)$ denotes the set of matchings on a collection consisting of $\al_i$ copies of each $i \in I$.
\end{definition}

\noindent When $|I| = 1$, $X\sim \calN(0,1)$, and $\Sig = 1$, then $\He_{(p)}(X\,;\, \Sig)$ equals the $p$th
Hermite polynomial.

When the $X_i$ are mean-zero Gaussian random variables and $\matSig$
is their covariance matrix, the Wick products satisfy the
(partial) orthogonality property that for each
finitely supported $\alpha, \beta\in\N^I$ with $\sum_i \alpha_i \neq \sum_i \beta_i$,
\[
    \E\left[\He_{\al}(\bX\,;\, \matSig) \He_{\beta}(\bX\,;\, \matSig)\right] = 0\,.
\]
In general, we have
\[ \E\left[\He_{\al}(\bX\,;\, \matSig) \He_{\beta}(\bX\,;\, \matSig)\right] = \sum_{M \in \PM(\alpha, \beta)} \prod_{uv \in M} \bm\Sigma[u, v]\,. \]
Since by the Wick lemma $\E\left[\prod_{i\in \alpha}X_i \cdot \prod_{j\in \beta} X_j\right]$ equals the same sum over all matchings of $\alpha \sqcup \beta$, the Wick products achieve a general ``partial orthogonalization'' that removes all terms from this covariance where any pairs within $\alpha$ or within $\beta$ are matched.

For each choice of $\matSig \in \R^{I \times I}_\sym$,
the Wick products are a basis for polynomials in the $X_i$.
Multiplication of polynomials gives an algebra structure to this space
which we call the {\em Wick algebra} of $\bX$.
Below is a combinatorial formula for multiplication in the Wick algebra.

\begin{proposition}[{\cite[Theorem 3.15]{Janson:GaussianHilbertSpaces}}]
\label{prop:wick}
    Let $I$ be an index set, $\bX = (X_i)_{i \in I}$ be formal variables, and $\matSig \in \R_\sym^{I \times I}$.
    Let $\al^1,\dots, \al^k \in \N^I$. Then:
    \begin{align*}
        \prod_{j=1}^k \He_{\al^j}(\bX; \matSig) = \sum_{M \in \calM( \al^1, \dots,  \al^k)} \prod_{uv \in M} \matSig[u, v] \He_{U(M)}(\bX\,;\, \matSig)\,,
    \end{align*}
     where $\al^j$ is a multiset of size $|\alpha^j|$ with $\al^j_{i}$ copies of each $i \in I$.
     Here $U(M) \in \N^I$ for $M$ a matching of $\al^1 \sqcup \cdots \sqcup \al^k$ counts the number of unmatched elements of each type.
\end{proposition}

\noindent In the special case where each
group $\alpha^j$ consists of a
single element, we obtain:

\begin{corollary}\label{lem:wick-recursion}
    For every $i_1, \ldots, i_k\in I$,
    \[
        \prod_{j=1}^k X_{i_j} = \sum_{M\in \calM(k)} \prod_{uv\in M} \bm \Sigma[i_u,i_v] \He_{U(M)}(\bm X\,;\,\bm \Sigma)\,.
    \]
\end{corollary}

\section{Diagrams and the \texorpdfstring{$w$- and $z$-Bases}{w- and z-Bases} of Polynomials}
\label{sec:diagrams}

All graphs considered in this paper are {\em multigraphs} (loops and multiedges are allowed) and will be
denoted by Greek letters ($\alpha, \beta, \gamma,\ldots$). 
We use the terms {\em graphs} and {\em diagrams}
interchangeably in this paper.
Given a diagram $\alpha$, we use $V(\alpha)$
to denote its vertex set and $E(\alpha)$ to denote
its edge set. We 
denote by $\alpha[S]$ the subgraph of $\alpha$
induced by $S\subseteq V(\alpha)$.
We count self-loops as contributing 2 to the degree of a vertex.

\subsection{Classes of diagrams}

Each diagram can have either $0$, $1$, or an ordered pair of $2$ special vertices called its \emph{root(s)}. With the exception of the class of graphs defined in~\Cref{def:open-cactus}, the roots of a graph can be arbitrary vertices (in particular, they might be equal if there are two of them).

\begin{notation}
    Let $\calA=\calA_0$ (resp.\ $\calA_1$ or $\calA_2$) be the set of all connected graphs with no 
    root (resp.\ $1$ root or $2$ roots). We also
    refer to such graphs as scalar (resp. vector or matrix) diagrams.
\end{notation}

Given $\alpha\in\calA$, an edge
$e\in E(\alpha)$ is a {\em bridge} of $\alpha$ if deleting $e$
would disconnect the graph. $\alpha\in\calA$ is {\em 2-edge-connected} if it contains no bridge. In general, $\alpha\in\calA$ can be decomposed into 
a tree of 2-edge-connected components connected by
bridges.

\begin{notation}
    Let $\calE=\calE_0\subseteq\calA$ (resp.\ $\calE_1\subseteq \calA_1$ or $\calE_2\subseteq \calA_2$) be the set of all 2-edge-connected
    scalar  (resp.\ vector or matrix) diagrams.
\end{notation}

Given $\alpha\in\calA$, a vertex $u\in V(\alpha)$ 
is an {\em articulation point} of $\alpha$ if removing $u$ and its incident edges
disconnects the graph. 
$\alpha$ is 
{\em 2-vertex-connected} if it has no articulation point.
Any $\alpha\in\calA$ decomposes into its 2-vertex-connected components
(blocks), which refine the 2-edge-connected components. The block-cut
graph (whose vertices are the articulation points and the blocks, with edges for
incidence) is a tree.

A connected graph is a {\em cactus} if every edge
lies on exactly one simple cycle.
Thus, cactuses are in a sense the minimal 2-edge-connected graphs.

\begin{notation}
    Let $\calC=\calC_0\subseteq \calA$ (resp. $\calC_1\subseteq \calA_1$) be
    the set of all scalar (resp. vector) cactus diagrams.
\end{notation}

\noindent For a cactus $\sigma$, we will denote
by $\cyc(\sigma)$ the set of (unrooted)
cycles of $\sigma$. 

Finally, as in~\Cref{def:treelike}, we will denote the treelike diagrams by $\calT_1$ and the treelike diagrams such that the root has degree 1 after deleting all hanging cactuses by $\calG_1$.

\subsection{Graph polynomials}
\label{sec:polynomial-definitions}

Each diagram represents different scalar-, vector-, or matrix-valued polynomials 
in a matrix input, depending on whether it is viewed
in the $w$-basis or the $z$-basis. In the following
definitions,
we fix $\bA\in \R_\sym^{n\times n}$, $\alpha$ to be a scalar, vector, or matrix diagram, and
$i,j\in [n]$.

\begin{definition}
    \label{def:w-basis}
    Define $w_\alpha(\bA)
    \in \R$, $\bw_\alpha(\bA)\in \R^n$, and $\bW_\alpha(\bA)\in \R^{n\times n}$ by
    \begin{alignat*}{2}
        w_\alpha(\bA) &= \sum_{\substack{\varphi\colon V(\alpha)\to [n]}} \prod_{\{u,v\}\in E(\alpha)} \bA[\varphi(u),\varphi(v)] \quad&& \text{if $\alpha$ is a scalar diagram,}\\[2mm]
        \bw_\alpha(\bA)[i] &= \sum_{\substack{\varphi\colon V(\alpha)\to [n]\\\varphi(r) = i}} \prod_{\{u,v\}\in E(\alpha)} \bA[\varphi(u),\varphi(v)] \quad&&\text{if $\alpha$ is a vector diagram with root $r$,}\\
        \bW_\alpha(\bA)[i,j] &= \sum_{\substack{\varphi\colon V(\alpha)\to [n]\\\varphi(r_1)=i,\varphi(r_2)=j}} \prod_{\{u,v\}\in E(\alpha)} \bA[\varphi(u),\varphi(v)] \quad&&\text{if $\alpha$ is a matrix diagram with roots $(r_1,r_2)$.}
    \end{alignat*}   
\end{definition}

\begin{definition}
    \label{def:z-basis}
    Define $z_\alpha(\bA)
    \in \R$, $\bz_\alpha(\bA)\in \R^n$, and $\bZ_\alpha(\bA)\in \R^{n\times n}$ by
    \begin{alignat*}{2}
        z_\alpha(\bA) &= \sum_{\substack{\varphi\colon V(\alpha)\hookrightarrow [n]}} \prod_{\{u,v\}\in E(\alpha)} \bA[\varphi(u),\varphi(v)] &&\text{if $\alpha$ is a scalar diagram,}\\[2mm]
        \bz_\alpha(\bA)[i] &= \sum_{\substack{\varphi\colon V(\alpha)\hookrightarrow [n]\\\varphi(r) = i}} \prod_{\{u,v\}\in E(\alpha)} \bA[\varphi(u),\varphi(v)] && \text{if $\alpha$ is a vector diagram with root $r$,}\\
        \bZ_\alpha(\bA)[i,j] &= \sum_{\substack{\varphi\colon V(\alpha)\hookrightarrow [n]\\\varphi(r_1)=i,\varphi(r_2)=j}} \prod_{\{u,v\}\in E(\alpha)} \bA[\varphi(u),\varphi(v)] \quad&& \text{if $\alpha$ is a matrix diagram with roots $(r_1,r_2)$.}
    \end{alignat*}    
\end{definition}

\noindent The only difference between the $w$- and $z$-bases
is the summation domain: \Cref{def:z-basis} sums over injective embeddings $\varphi$, whereas \Cref{def:w-basis} sums over all
embeddings.

Finally, we define two extensions of~\Cref{def:w-basis}
that we will need in the proofs.
The following allows us to use a different matrix on
each edge of the graph:

\begin{definition}
    \label{def:non-uniform-label}
    Let $\alpha$ be a matrix diagram with roots $(r_1,r_2)$ and $\bm \calA = (\bmA_e)_{e\in E(\alpha)}$ be
    such that $\bmA_e\in \R_\sym^{n\times n}$ for
    all $e\in E(\alpha)$. Define $\bW_\alpha(\bm \calA)\in \R^{n\times n}$ by
    \[
        \bW_\alpha(\bm \calA)[i,j] = \sum_{\substack{\varphi\colon V(\alpha)\to [n]\\\varphi(r_1)=i,\varphi(r_2)=j}} \prod_{e=\{u,v\}\in E(\alpha)} \bA_e[\varphi(u),\varphi(v)]\,.
    \]
\end{definition}

The following is an intermediate quantity
between~\Cref{def:w-basis} and~\Cref{def:z-basis} which only restricts the sum
over injective labelings on two vertices:

\begin{definition}
    Let $\bA\in \R_\sym^{n\times n}$, $\alpha$ be a scalar/vector/matrix diagram,
    $i,j\in [n]$, and $s,t\in V(\alpha)$. Define $w_\alpha^{s\neq t}\in \R$, $\bw_\alpha^{s\neq t}\in \R^n$, and $\bW_\alpha^{s\neq t}\in \R^{n\times n}$ by
    \begin{alignat*}{2}
        w_\alpha^{s\neq t}(\bA) &= \sum_{\substack{\varphi\colon V(\alpha)\to [n]\\\varphi(s)\neq \varphi(t)}} \prod_{\{u,v\}\in E(\alpha)} \bA[\varphi(u),\varphi(v)] && \text{if $\alpha$ is a scalar diagram,}\\
        \bw^{s\neq t}_\alpha(\bA)[i] &= \sum_{\substack{\varphi\colon V(\alpha)\to [n]\\\varphi(s)\neq \varphi(t)\\\varphi(r) = i}} \prod_{\{u,v\}\in E(\alpha)} \bA[\varphi(u),\varphi(v)] && \text{if $\alpha$ is a vector diagram with root $r$,}\\
        \bW^{s\neq t}_\alpha(\bA)[i,j] &= \sum_{\substack{\varphi\colon V(\alpha)\to [n]\\\varphi(s)\neq \varphi(t)\\\varphi(r_1)=i\\\varphi(r_2)=j}} \prod_{\{u,v\}\in E(\alpha)} \bA[\varphi(u),\varphi(v)] \quad && \text{if $\alpha$ is a matrix diagram with roots $(r_1, r_2)$.}
    \end{alignat*}
\end{definition}

\subsection{Partitions, change of basis, and M{\"o}bius inversion}
\label{sec:mobius}

While $(z_\alpha(\bA))_{\alpha\in \cA}$ and
$(w_\alpha(\bA))_{\alpha\in \cA}$ span the
same space of $S_n$-invariant polynomials in the entries of $\bm A$, some
properties are better expressed in one basis
than the other. 
Here we take a closer look at these bases and derive
change-of-basis formulas.

Given a set $S$, let $\calP(S)$ denote the set of all \emph{partitions} of $S$, sets of non-empty disjoint subsets of $S$ whose union is all of $S$.
We call the parts of a partition \emph{blocks}.
Each block is a set, and $P$ is the set of blocks, so we denote the blocks by $b \in P$.

For a (scalar, vector, or matrix) diagram $\al$ and a partition $P \in \calP(V(\al))$,
we define a new diagram $\al_P$
by identifying the vertices within each block of $P$ into a single vertex.
The vertices of $\al_P$ may thus be identified with the blocks of $P$.
$\al_P$ retains all edges of $\al$, which may become multiedges or self-loops.
The status of being one of the (0, 1, or 2) roots of $\al$ is inherited by the block containing that root. 

To change from the $w$- to the $z$-basis, we then simply sum over all partitions:

\begin{claim}
    \label{claim:wtoz}
    For all (scalar, vector, or matrix) diagrams $\alpha$, 
    \[
        w_{\alpha}(\bA) = \sum_{P \in \calP(V(\alpha))} z_{\alpha_P}(\bA)\,.
    \]
\end{claim}

Define the relation $\al \psdleq \beta$ on scalar diagrams if there exists a partition $P \in \calP(V(\beta))$
such that $\al = \beta_P$.
It is easy to check that this relation gives a partial ordering, inherited from the standard partial ordering on partitions.
We write $\alpha\prec \beta$ as a shorthand for
$\alpha\preceq \beta$ and $\alpha\neq \beta$.

\begin{lemma}
    \label{claim:mobius}
    There exist $(c_{\alpha, \beta})_{\alpha, \beta \in \calA}$ and $(c'_{\alpha, \beta})_{\alpha,\beta\in\calA}$ not depending on $n$ such that $c_{\alpha, \beta} \in \N$, $c'_{\al,\beta} \in \Z$ and for any $\alpha,\beta\in\calA$, 
    \[
        w_\beta(\bA) = \sum_{\al \preceq \beta} c_{\alpha, \beta} z_\alpha(\bA)\,,\qquad 
        z_{\beta}(\bA) = \sum_{\alpha \preceq \beta} c'_{\alpha,\beta} w_{\alpha}(\bA)\,.
    \]
\end{lemma}
\begin{proof}
    The coefficients in the left equation count symmetries in \cref{claim:wtoz}, i.e.,
    $c_{\al,\beta}$ equals the number of ways to choose a partition $P \in \calP(V(\beta))$
    such that $\beta_P$ is isomorphic to $\alpha$.
    Reciprocally, since $\preceq$ is a partial ordering, this
    transformation can be inverted using M{\"obius} inversion~\cite{Rota-1964-Foundations} on this poset.
    Although an explicit formula for $c'_{\al, \beta}$
    is available in terms of the
    combinatorial structure of the graphs, we will not need it in this paper.
\end{proof}


\subsection{The example of cycles: Moments versus free cumulants}
\label{sec:moments-free-cumulants}

The difference
between the $w$- and $z$-bases is illustrated nicely by the
special case of the diagrams $\sigma_q$ which are cycles of length $q\ge 1$.
In this case, $\frac 1n w_{\sigma_q}(\bA)$ and $\frac 1n z_{\sigma_q}(\bA)$ are versions of the limiting spectral moments and free cumulants, respectively, for finite-dimensional matrices.


Let $\calP(q)$ denote the set of partitions of $\{1,2,\dots, q\}$
and let $\textnormal{NC}(q)$ denote the subset of non-crossing partitions (partitions such that there does not exist $i < j < k < \el$ with $i,k$ in the same block and $j, \el$ in the same block, different from the one $i,k$ are in).
It is convenient to view these as partitions of the vertices of the $q$-cycle
so that the term \emph{non-crossing} may be interpreted visually: in a non-crossing partition, the blocks do not intersect one another when drawn as ``blobs'' inside the cycle.

In the $w$-basis, we have
\begin{align}
    \frac 1n w_{\sigma_q}(\bA) = \frac 1n \sum_{i_1,\ldots,i_q=1}^n \bA[i_1, i_2] \bA[i_2, i_3] \ldots \bA[i_q,i_1] = \frac 1n \Tr(\bA^q) = \frac 1n \sum_{i=1}^n \lambda_i(\bA)^q\,.\label{eq:empiricalDist}
\end{align}
Suppose that the expression in~\cref{eq:empiricalDist} converges as
$n\to\infty$ to the $q$th moment
$m_q\in \R$ of a limiting spectral distribution, $m_q = \int \lambda^q\,\d \mu(\lambda)$.

The free cumulants
are defined from the moments by a formula similar to
the classical cumulants vis-{\`a}-vis the moments of a random variable:

\begin{definition}[Free cumulant]\label{def:free-cumulant}    
    The free cumulants $(\kappa_q)_{q\ge 1}$ corresponding to $(m_q)_{q\ge 1}$ are defined implicitly by:
\begin{equation}\label{eq:free-cumulant}
    m_q = \sum_{\sig \in \textnormal{NC}(q)} \prod_{b \in \sig} \kappa_{|b|}\,.
\end{equation}
\end{definition}
\noindent The $\kappa_q$ can be computed explicitly in terms of the $m_q$ by applying M{\"o}bius inversion to~\cref{eq:free-cumulant}; see~\cref{eq:free-cumulant-explicit}.

Analogous to \cref{eq:empiricalDist} which is in the $w$-basis,
it appears to be folklore\footnote{This is for example explicitly stated in \cite[Theorem 1 and Appendix D.1]{MFCKMZ-2019-PlefkaExpansionOrthogonalIsing}.} that if $\bA$ is drawn
from an orthogonally invariant matrix ensemble
with free
cumulants $(\kappa_q)_{q\ge 1}$, then
\begin{align}
    \frac 1n \E z_{\sigma_q}(\bA) \underset{n\to\infty}{\longrightarrow} \kappa_q\,.\label{eq:freeCumulantRelation}
\end{align}
The quantity $\frac 1n z_{\sigma_q}(\bA)$ has also been called the $q$th \emph{injective trace} of $\bA$.
Below in \cref{lem:diagonal-w-z}, we prove \cref{eq:freeCumulantRelation} using a change of basis from $w$ to $z$.

For example, below are the parameters $m_q$ and $\kappa_q$
for the \goe and the \rom, whose limiting empirical spectral distribution
are the Wigner semicircle distribution and the Rademacher distribution, respectively.

\begin{claim}\label{claim:rom-cumulants}
    Let $\textnormal{Cat}(k) \colonequals \frac{1}{k + 1}\binom{2k}{k}$ be the $k$th Catalan number.
    For the \goe, the limiting spectral moments and free cumulants are:
    \begin{align*}
        m_q = \left\{\begin{array}{ll}
            \textnormal{Cat}(q/2) & \textnormal{if $q$\text{ is even}}\\
            0 & \textnormal{if $q$\text{ is odd}}
        \end{array}
        \right\}, \qquad 
        \kappa_q = \left\{\begin{array}{ll}
            1 & \textnormal{if $q = 2$}\\
            0 & \textnormal{otherwise}
        \end{array}\right\}.
    \end{align*}        
    For the \rom, the limiting spectral moments and free cumulants are:
    \begin{align}
    m_q = \left\{\begin{array}{ll}
        1 & \textnormal{ if $q$ is even}\\
        0 & \textnormal{ if $q$ is odd}
    \end{array}\right\}, \qquad
    \kappa_q = \left\{\begin{array}{ll}
        (-1)^{q/2-1} \textnormal{Cat}(q/2 - 1) & \textnormal{ if $q$ is even}\\
        0 & \textnormal{ if $q$ is odd}
    \end{array}\right\}.\label{eq:freeCumulantROM}
\end{align}
\end{claim}


\subsection{Solving equations in the traffic distribution}

The traffic distribution is defined as the limiting values of all $w$-basis polynomials, but we show now how it can be derived from various combinations of limits of $w$- and $z$-basis polynomials.
In our other arguments, we will also find it convenient to describe
the traffic distribution of sequences of matrices (random or deterministic)
using the two bases simultaneously.
While~\Cref{claim:mobius} shows
that we could in principle express all these results
in a single basis, this would involve
precisely tracking 
very complicated combinatorial coefficients
(in fact, this was a major technical obstacle in
previous diagrammatic analyses of AMP).

As we have discussed, when a matrix satisfies the strong cactus property,
its traffic distribution is determined by its values on the cactus diagrams (equivalently, by the diagonal distribution),
and when it satisfies the factorizing strong cactus property,
its traffic distribution is determined by the spectral distribution.
We show that one can use either the $w$-basis or $z$-basis for these determinations.

\begin{lemma}
    \label{lem:diagonal-w-z}
    Suppose that $\bA = \bA^{(n)}$ satisfies the
    weak cactus property, i.e., for all $\alpha\in \calE\setminus \calC$,
    \[
        \frac 1n \E_{\bmA} z_\alpha(\bmA)\underset{n\to\infty}{\longrightarrow} 0\,.
    \]
    Then the following are equivalent:
    \begin{enumerate}[(i)]
        \item For all $\sigma \in \calC$ there exists $m_\sigma \in \R$ such that $\frac 1n \E_\bA w_\sigma(\bA) \underset{n\to\infty}{\longrightarrow} m_\sigma$.
        \item For all $\sigma \in \calC$ there exists $k_\sigma \in \R$
        such that $\frac 1n \E_\bA z_\sigma(\bA) \underset{n\to\infty}{\longrightarrow} k_\sigma$.
    \end{enumerate}
    Furthermore, when they exist, $(m_\sigma)_{\sigma \in \calC}$ and $(k_\sigma)_{\sigma \in \cC}$ determine each other.
    The following are also equivalent:
    \begin{enumerate}[(i)]
        \item There exist real numbers $(m_q)_{q \in \N}$ such that for all $\sig \in \calC$, $\frac 1n \E_\bA w_\sig(\bA) \underset{n\to\infty}{\longrightarrow} \prod_{\rho \in \cyc(\sig)} m_{|\rho|}$.
        \item There exist real numbers $(\kappa_q)_{q \in \N}$ such that for all $\sig \in \calC$, $\frac 1n \E_\bA z_\sig(\bA) \underset{n \to \infty}{\longrightarrow} \prod_{\rho \in \cyc(\sig)} \kappa_{|\rho|}$.
    \end{enumerate}
    Furthermore, when they exist, $(m_q)_{q \in \N}$ and $(\kappa_q)_{q \in \N}$ are related by \cref{eq:free-cumulant}.
\end{lemma}

We use the
following observation which will be used repeatedly in \cref{sec:universality}:

\begin{lemma}\label{lem:contract-2ec}
    If $\al \in \calE$ and $\beta \psdleq \al$, then $\beta\in\calE$.
\end{lemma}

\begin{proof}[Proof of~\Cref{lem:contract-2ec}]
    By Menger's theorem, a graph is 2-edge-connected 
    if and only if there exist two edge-disjoint 
    paths between every pair of distinct vertices.
    These paths are maintained when $\al$ is contracted into $\beta$.
\end{proof}

\begin{proof}[Proof of~\Cref{lem:diagonal-w-z}]
    {\em (ii)} $\implies$ {\em (i)}. Using \Cref{claim:wtoz},
    \[
        \frac 1n \E_{\bA} w_\sigma(\bA) = \frac 1n \sum_{\beta \psdleq \sigma} c_{\beta, \sigma} \E_{\bA} z_\beta(\bA)\,.
    \]
    Every diagram $\beta \psdleq \sigma$ remains 2-edge-connected by \cref{lem:contract-2ec}.
    There are only finitely many terms in the
    sum, so we can directly
    take the $n\to\infty$ limit
    and use the assumptions
    to obtain that $\frac 1n \E_\bA w_\sigma(\bA)$ converges to $\sum_{\beta \psdleq \sig} c_{\beta, \sig} k_\beta$.

    Note that by the weak cactus property,
    the only asymptotically nonzero $\beta \preceq \sig$ are when $\beta$ is a cactus.
    Assuming furthermore that $k_\beta = \prod_{\rho \in \cyc(\beta)} \kappa_{|\rho|}$ factors over the cycles of each cactus $\beta$ we will derive the second part of the lemma.

    Using the more specific result of \Cref{claim:wtoz}, we have
    \begin{align*}
        \lim_{n \to \infty} \frac{1}{n} \E_{\bA} w_{\sigma}(\bA)
        &= \sum_{P \in \calP(V(\sigma))} \lim_{n \to \infty} \frac{1}{n} \E_{\bA} z_{\sigma_P}(\bA)
        \intertext{Since $\bA$ has the weak cactus property and $\sigma$ is a cactus, only the terms where $\sigma_P$ is a cactus contribute. These are precisely the terms where $P$ restricted to each cycle of $\sigma$ is non-crossing. Given $P_{\rho} \in \NC(V(\rho))$ for each $\rho \in \cyc(\sigma)$, let us write $P(P_{\rho}: \rho \in \cyc(\sigma))$ for the partition obtained by composing these partitions of each cycle, and let us write, following our previous notation, $\cyc(\rho_{P_{\rho}})$ for the set of cycles created when the single cycle $\rho$ is contracted according to $P_{\rho}$. Then, we have}
        &= \sum_{\substack{P_{\rho} \in \NC(V(\rho)) \\ \text{for each } \rho \in \cyc(\sigma)}} \lim_{n \to \infty} \frac{1}{n} \E_{\bA} z_{\sigma_{P(P_{\rho}: \rho \in \cyc(\sigma))}}(\bA) \\
        &= \sum_{\substack{P_{\rho} \in \NC(V(\rho)) \\ \text{for each } \rho \in \cyc(\sigma)}} \prod_{\rho \in \cyc(\sigma)} \prod_{\pi \in \cyc(\rho_{P_{\rho}})} \kappa_{|\pi|} \\
        &= \prod_{\rho \in \cyc(\sigma)} \left(\sum_{P \in \NC(V(\rho))} \prod_{\pi \in \cyc(\rho_P)} \kappa_{|\pi|}\right) \\
        &= \prod_{\rho \in \cyc(\sigma)} m_{|\rho|}.
    \end{align*}
    Thus we have the claimed factorization.
    Further, the coefficients $m_q$ and $\kappa_q$ indeed have the relation between moments and free cumulants from~\cref{eq:free-cumulant}:
    \[
        m_q = \sum_{\sig \in \NC(q)} \prod_{b \in \sig} \kappa_{|b|}\,.
    \]
    
    \noindent {\em (i) $\implies$ (ii)}. This direction uses a recursive change of basis technique that will be very useful in~\Cref{sec:universality}.
    Using~\Cref{claim:mobius} in both directions, we get
    \begin{align*}
        \frac 1n \E_{\bmA} z_\sigma(\bmA) &= \frac 1n \sum_{\substack{\beta\preceq \sigma\\\beta\in\calC}} c'_{\beta,\sigma} \E_{\bmA} w_\beta(\bmA) + \frac 1n \sum_{\substack{\beta\prec \sigma\\\beta\in\calE\setminus\calC}} c'_{\beta,\sigma} \E_{\bmA} w_\beta(\bmA)\\
        &= \frac 1n \sum_{\substack{\beta\preceq \sigma\\\beta\in\calC}} c'_{\beta,\sigma} \E_{\bmA} w_\beta(\bmA) + \frac 1n \sum_{\substack{\beta\prec \sigma\\\beta\in\calE\setminus\calC}} c'_{\beta,\sigma} 
        \sum_{\alpha\preceq \beta} c_{\alpha,\beta} \E_{\bmA} z_\alpha(\bmA)\\
        &= \frac 1n \sum_{\substack{\beta\preceq \sigma\\\beta\in\calC}} c'_{\beta,\sigma} \E_{\bmA} w_\beta(\bmA) +  \frac 1n  \sum_{\alpha\prec \sigma}\left(\sum_{\substack{\beta\in\calE\setminus\calC\\ \alpha \preceq \beta \prec \sigma}} c'_{\beta,\sigma} c_{\alpha,\beta}\right) \E_{\bmA} z_\alpha(\bmA)
    \end{align*}
    Note that every diagram in this expansion
    remains 2-edge-connected by~\Cref{lem:contract-2ec}.
    
    Every contraction identifying
    a non-empty subset of vertices decreases
    the number of vertices in the graph, and
     the $w$ and $z$ bases coincide for 1-vertex
    graphs. Therefore, we can
    apply the same steps inductively on terms
    for which $\alpha\in\calC$ to finally obtain
    \[
        \frac 1n \E_{\bmA} z_\sigma(\bmA) = \frac 1n \sum_{\substack{\beta\preceq \sigma\\\beta\in \calC}} c''_{\beta,\sigma} \E_{\bmA} w_\beta(\bmA) + \frac 1n \sum_{\substack{\beta\preceq \sigma\\\beta\in \calE\setminus \calC}} c''_{\beta,\sigma} \E_{\bmA} z_\beta(\bmA)\,.
    \]
    for some coefficients $\{c''_{\alpha,\beta}\}$ independent of $n$.
    Take the $n\to\infty$ limit
    to obtain
    \[
        \lim_{n\to\infty} \frac 1n \E_{\bmA} z_\sigma(\bmA) = \sum_{\substack{\beta\preceq \sigma\\\beta\in \calC}} c''_{\beta,\sigma} m_\beta\,,
    \]
    which finishes the proof of the first equivalence.
    Assuming furthermore that $m_\beta$ factors over the cycles of each cactus $\beta$,
    then $\frac 1n \E_\bA z_\sig(\bA)$
    also asymptotically factors over its cycles: $\frac 1n \E_\bA z_\sig(\bA) \longrightarrow \prod_{\rho \in \cyc(\sig)} \kappa_{|\rho|}$ for some numbers $\kappa_q$.
    This is because the cactuses $\beta \psdleq \sig$ still only arise by contracting a separate non-crossing partition for each cycle of $\sig$, and so we can perform the above recursive analysis separately inside each cycle.
\end{proof}

The following lemma shows that the properties
of graph polynomials
we will establish for delocalized deterministic matrices
in~\Cref{sec:universality} characterize their traffic distribution.
We emphasize our use of a combination of assumptions on limits of the $w$- and $z$-bases that makes this formulation convenient.

\begin{lemma}
    \label{lem:eqDiffMobius}
    Suppose that $\bA = \bA^{(n)}$ satisfies:
    \begin{enumerate}
        \item The weak cactus property, i.e., that for all $\alpha\in \calE\setminus \calC$, $\frac 1n \E_{\bmA} z_\alpha(\bmA)\underset{n\to\infty}{\longrightarrow} 0$.
        \item For all $\alpha\in \calA\setminus \calE$, $\frac 1n \E_{\bmA} w_\alpha(\bmA)\underset{n\to\infty}{\longrightarrow} 0$.
        \item For all $\sigma\in \calC$,
        there exists $m_\sigma\in \R $ such that
        $\frac 1n \E_{\bmA} w_\sigma(\bmA)\underset{n\to\infty}{\longrightarrow} m_\sigma$.
    \end{enumerate}
    Then the traffic distribution of $\bmA$
    exists
    and only depends on $\{m_\sigma : \sigma\in\calC\}$.
\end{lemma}

\begin{proof}
    We want to show that for every $\alpha\in \calA$, $ \lim_{n\to\infty} \frac 1n \E_{\bmA} w_\alpha(\bmA)$ exists and
    only depends on $\{m_\sigma : \sigma\in\calC\}$. By assumption, it suffices to prove it for $\alpha\in \calE\setminus \calC$. By~\Cref{claim:mobius},
    \[
        \frac 1n \E_{\bmA} w_\alpha(\bmA) = \frac 1n \sum_{\beta\preceq \alpha} c_{\beta,\alpha} \E_{\bmA} z_\beta(\bmA)\,.
    \]
    By~\Cref{lem:contract-2ec}, every $\beta$
    in the support of the sum is 2-edge-connected. If $\beta\in\calC$, then the value of $\lim_{n\to\infty} \frac 1n  \E_{\bmA} z_\beta(\bmA)$ exists and only
    depends on $\{m_\sigma : \sigma\in\calC\}$
    by~\Cref{lem:diagonal-w-z}. Otherwise,
    $\beta\in\calE\setminus \calC$, and $\lim_{n\to\infty} \frac 1n  \E_{\bmA} z_\beta(\bmA) = 0$ by assumption. This
    implies that $\lim_{n\to\infty}\frac 1n \E_{\bmA} w_\alpha(\bmA)$ exists and 
    only depends on $\{m_\sigma : \sigma\in\calC\}$, which concludes the proof.
\end{proof}
\noindent
Note that, more generally, by~\Cref{lem:diagonal-w-z}, the same statement will hold with Condition~3 of~\Cref{lem:eqDiffMobius} taken in terms of either the $w$- or $z$-basis.


\subsection{Products and concentration of traffic observables}

Recall that the traffic distribution specifies the limits of $\frac 1n \E_\bA w_\al(\bA)$
for all $\al \in \calA$.
In all of the random matrix models we consider,
these expectations are highly concentrated.
We say that {\em the traffic distribution concentrates for $\bA$} if the
following property holds, studied in~\cite{male2020traffic}.

\begin{definition}\label{def:traffic-concentration}
    Let $\bA = \bA^{(n)} \in \R^{n \times n}_\sym$ and assume that the traffic distribution of $\bA$ exists. We say that the traffic distribution concentrates for $\bA$
    if for all $k \geq 2$ and $\al_1, \dots, \al_k \in \calA$,
    \[
        \lim_{n \to \infty} \E_{\bA} \left[ \prod_{j=1}^k \frac 1n w_{\alpha_j}(\bA)\right] = \prod_{j=1}^k \lim_{n \to \infty} \frac 1n \E_{\bA} w_{\alpha_j}(\bA)\,.
    \]
\end{definition}
\noindent
The case $k = 2$ and $\alpha_1 = \alpha_2 = \alpha$ of the definition specializes to the statement:
\begin{lemma}\label{lem:traffic-l2}
    Let $\bA = \bA^{(n)} \in \R^{n \times n}_{\sym}$ have traffic distribution $\calD$. 
    If the traffic distribution concentrates for $\bA$, then $\frac 1n \E_{\bA} w_\al(\bA)$ converges to $\calD(\al)$ in $L^2$.
\end{lemma}
The full condition may be viewed as a strengthening of this straightforward notion of concentration.
We note that the product of several $w$-basis polynomials is equivalent to taking the disjoint union of their diagrams:
\[ w_{\alpha_1}(\bA) \cdots w_{\alpha_k}(\bA) = w_{\alpha_1 \sqcup \cdots \sqcup \alpha_k}(\bA).\]
Therefore, \cref{def:traffic-concentration} says that the values of disconnected diagrams asymptotically factor over the components.
This justifies defining
$\cA$ and the traffic distribution to include only connected diagrams.
The following shows that concentration may equally well be considered in the $z$-basis.

\begin{lemma}[{\cite[Lemma 2.9]{male2020traffic}}] \label{lem:concentration-z}
    Let $\bA = \bA^{(n)} \in \R^{n \times n}_\sym$ and assume that the traffic distribution of $\bA$ exists. The traffic distribution concentrates for $\bA$ if and only if, for all $k \geq 2$ and $\al_1, \dots, \al_k \in \calA$,
    \[
        \lim_{n \to \infty} \E_{\bA} \left[ \prod_{j=1}^k \frac 1n z_{\alpha_j}(\bA)\right] = \prod_{j=1}^k \lim_{n \to \infty} \frac 1n \E_{\bA} z_{\alpha_j}(\bA)\,.
    \]
\end{lemma}


For vector diagrams, the componentwise or Hadamard product is
\[
    \bmw_{\al_1}(\bA) \cdots \bmw_{\al_k}(\bA) = \bw_{\al_1 \oplus \cdots \oplus \al_k}(\bA)\,,
\]
where $\al_1 \oplus \cdots \oplus \al_k$ is the diagram formed by taking the disjoint union of $\al_1$ through $\al_k$ and then identifying the roots together into a single root.
We sometimes refer to this operation as \emph{grafting} $\al_1, \dots, \al_k$ at the root.